\subsubsection{複素数と複素平面}
  まず,虚数を定義する.
  虚数とは,2\,乗して\,$-1$になる数のことで,多くの場面で\,$i$を用いて
  \[
    \sqrt{-1}=i
  \]
  と定義される.
  複素数は
  \[
    z=a+bi\qquad (a,b\in\mathbb{R},\ i^2=-1)
  \]
  と表される数である.ここで,\,$a$を実部,\,$b$を虚部といい,
  \[
    \operatorname{Re}(z)=a,\qquad \operatorname{Im}(z)=b
  \]
  と表す.\,$\operatorname{Re}$は実数(Real number)に由来し,\,$\operatorname{Im}$は虚数(Imaginary number)に由来する.

  複素数\,$z=a+bi$は,複素平面上の点$(a,b)$と対応させることができる.横軸を実軸,縦軸を虚軸とすると,\,$z=a+bi$は点$(a,b)$を表す.

  特に,複素数\,$z=a+bi$の絶対値は
  \[
    |z|=\sqrt{a^2+b^2}
  \]
  であり,原点から対応する点までの距離を表す.また,\,$z\neq0$に対して
  \[
    z=|z|(\cos\theta+i\sin\theta)
  \]
  と表すことができる.ここで,$\theta$は正の実軸から測った偏角である.この表示を極形式という.

  オイラーの公式
  \[
    e^{i\theta}=\cos\theta+i\sin\theta
  \]
  を用いれば,
  \[
    z=|z|e^{i\theta}
  \]
  と簡潔に表すことができる.この表示は,複素数の積を考えるときに特に便利である.実際,
  \[
    z_1=r_1e^{i\theta_1},\qquad z_2=r_2e^{i\theta_2}
  \]
  ならば,
  \[
    z_1z_2=r_1r_2e^{i(\theta_1+\theta_2)}
  \]
  となる.つまり,複素数を掛けることは,絶対値を掛け合わせると同時に,偏角を加えることに対応する.

  また,複素数\,$z=a+bi$に対して,
  \[
    \overline{z}=a-bi
  \]
  をその共役複素数という.このとき,
  \[
    z\overline{z}=a^2+b^2=|z|^2
  \]
  が成り立つ.

\subsubsection{共役な複素数}
  実係数の2\,次方程式
  \[
    x^2+px+q=0\qquad (p,q\in\mathbb{R})
  \]
  が虚数解をもつ場合を考える.判別式が負であれば,
  \[
    p^2-4q<0
  \]
  であり,解は
  \[
    x=\frac{-p\pm\sqrt{p^2-4q}}{2}
    =\frac{-p}{2}\pm\frac{\sqrt{4q-p^2}}{2}i
  \]
  と書ける.したがって,解は
  \[
    \alpha=a+bi,\qquad \beta=a-bi
  \]
  という共役な複素数の組になる.

  これは,実係数の方程式が虚数解をもつときに重要な性質である.
  すなわち,実係数の方程式の非実数解は,共役な複素数として現れる.